\section{总结与展望}
\label{sec:conclusion}

\subsection{工作总结}

本文在 ChnSentiCorp 数据集上对 SVM + TF-IDF、Bi-LSTM、BERT-base 微调、Qwen2.5-7B 零样本和 Qwen2.5-7B LoRA 微调这五类模型进行了横向对比。实验在统一环境与超参数下完成，引入了早停机制确保公平。主要结论：

\textbf{（1）预训练是精度提升的最大推动力。} 从 Bi-LSTM（F1=0.9003）到 BERT（F1=0.9521），预训练范式带来了 5.18 个百分点的 F1 增益，占实验全程性能提升的大部分。在大规模无标注语料上积累的通用语义知识，是下游小样本任务的质量基石。

\textbf{（2）LoRA 高效释放了大模型的潜力。} Qwen2.5-7B 经 LoRA 微调后，F1 从 0.8985 升至 0.9570（+5.85\%），在全部四项指标上达到最优。以极低的参数量实现了显著的性能跃迁，验证了 PEFT 技术路线的有效性。

\textbf{（3）精度-效率的权衡清晰可见。} Qwen-LoRA 以 7.66 倍于 BERT 的训练成本换取 0.49 个百分点的 F1 提升——边际收益递减明显。在约 1 万条标注数据的条件下，BERT 是精度与效率的最佳平衡点。

\textbf{（4）传统方法和零样本在特定场景仍有意义。} SVM（F1=0.8944, 训练 0.32s）的计算经济性无可替代；Qwen-Zero（F1=0.8985）零训练成本给出有监督浅层模型级别的精度，适合标注完全不可得的场景。

\subsection{工程实践启示}

本研究的发现对实际工程选型具有以下指导意义：

\textbf{（1）数据规模决定模型上限的选择余量。} 在标注量约 1 万条的规模上，BERT 已经接近性能天花板。如果预算允许人工标注更多数据（如 5 万至 10 万条），Qwen-LoRA 与 BERT 之间的精度差距有可能进一步拉开——因为大模型在数据量增加时通常表现出更好的可扩展性。考虑到大模型的部署和维护成本，在数据量不多时盲目追求大模型可能并不划算。

\textbf{（2）业务场景决定指标偏重。} 舆情监控场景更关注精确率——误将中性或正向内容标为负面可能导致错误的商业决策，因此 BERT 和 Qwen-LoRA 这类 Precision 更高的模型更合适。内容审核场景更看重召回率——漏掉真正负面的内容可能带来合规风险，对模型的要求是"宁可错判，不可漏过"。自动化客服场景则对推理延迟有严格要求：SVM 或 Bi-LSTM 可以在 CPU 上以毫秒级响应，而 Qwen 系列在无 GPU 加速时难以满足实时性。

\textbf{（3）迭代周期决定技术门槛。} SVM 和 Bi-LSTM 不需要 GPU，训练和调参可以在普通 PC 上完成，适合快速原型开发和频繁迭代。BERT 和 Qwen-LoRA 需要 GPU 训练环境，但 BERT 的微调效率高（本实验中约 3 分钟），适合每周甚至每天的模型更新。Qwen-LoRA 的训练时间约 24 分钟，虽然可控，但在需要高频迭代的场景中会累积可观的等待时间。

\textbf{（4）零样本是标注冷启动的有效手段。} 在没有标注数据的初期，可用 Qwen-Zero 快速建立基线，收集线上反馈后再用 LoRA 微调持续提升效果——这种"零样本启动 + 有监督迭代"的路径在标注资源有限的项目中具有现实意义。

\subsection{研究局限}

\textbf{（1）数据集限制。} 实验限于 ChnSentiCorp 酒店评论语料，情感表达比较直接。结论能否推广到金融、医疗等垂直领域，或长文本、多分类任务，需要更多数据验证。

\textbf{（2）LoRA 配置不够丰富。} 目前只在 $r=16$、适配 Query 和 Value 的配置下实验。更优的秩、更广的模块选择（如增加 Key 和 Output 投影）、更多的标注数据，都有可能进一步拉开 Qwen-LoRA 与 BERT 的差距。

\textbf{（3）推理评估不够细。} 只统计了端到端推理耗时，没有区分分词、KV-Cache 和矩阵运算各自的占比，也没有测试批量推理（Batch Inference）的吞吐表现。

\textbf{（4）零样本 Prompt 未充分探索。} 只用了固定的一套提示词，没有尝试不同 Prompt 模板对 Qwen-Zero 精度的潜在提升。

\subsection{未来展望}

\textbf{（1）跨领域验证。} 在电商、社交媒体、金融评论、医疗问诊等多领域语料上做统一对比，检验各模型在不同文本分布、不同情感粒度和不同语言风格下的表现稳定性。特别值得关注的是：当文本复杂度上升（如长文档、多段落），各模型之间的差距是会扩大还是缩小。

\textbf{（2）LoRA 超参数全面搜索。} 对秩 $r$（8, 16, 32, 64）、适配模块组合（Q/K/V/O 的各种排列）、学习率和 Dropout 率做网格搜索，找出中文情感分类任务上的最优 LoRA 配置。此外还可以探索 LoRA+、AdaLoRA 等改进变体，看其在当前任务上是否有额外收益。

\textbf{（3）Prompt 工程与少样本学习。} 系统比较多组 Prompt 模板（如角色扮演型、思维链型、标签条件型等）下的零样本和少样本表现，考察 In-Context Learning 能否在不动模型权重的前提下显著缩小与大模型微调的差距。若能通过精心设计的 Few-shot Prompt 使 Qwen-Zero 的 F1 提升到 0.93 以上，则"零训练成本 + 高精度"的路径将更具工程吸引力。

\textbf{（4）推理加速与轻量化部署。} 尝试 4-bit/8-bit 量化、FlashAttention 加速和 vLLM 推理框架，评估 Qwen-7B 在保持精度的前提下能否将推理延迟压缩到 BERT 的 2 倍以内。若量化后精度损失可控（F1 降幅 < 0.01），则大模型的部署门槛将大幅降低。

\textbf{（5）多任务与少样本泛化的统一评估。} 将当前二分类实验框架扩展为统一的多任务评测基准，涵盖情感二分类、细粒度情感多分类（1-5 星）、方面级情感分析（Aspect-Based Sentiment Analysis）和情感原因抽取，更全面地衡量各模型在中文情感理解全链条上的能力差异。
